% s-197-p197-m001; original PDF 197, 198, 199
\Aphorism{s-197-p197-m001}{283}
Im Leben sind wir schon so klug geworden, nicht mehr „aufs Wort“ zu glauben, was man uns sagt; aus Misstrauen gegen das Wort an sich, nicht gerade gegen den Wortmacher selbst. Die Gründe, nach denen Einer gehandelt haben will, sind nicht die, nach denen er gehandelt hat; bei redlichster Bemühung, Thatbestände mit Worten zu decken, ist eine Congruenz eben doch nicht zu erzielen. Das wissen wir im Leben, und in der Kunst wollen wir thun, als wüssten wir es nicht? Was wir unserem Nachbar nicht glauben würden, wollen wir in aller Naivität dem verkleideten Acteur glauben oder Dem, der unsichtbar hinter ihm steht? Wir wollen so primitiv von Ursache und Wirkung denken, ein Geschehniss für motivirt zu halten, wenn es explicirt worden ist? Wo bleibt bei dieser Selbsttäuschung die Psychologie, in deren Namen man sich täuschen lässt? — Man kann viel Kunstgeschwätz sparen, wenn man sich die Analogiefrage stellt: wollen wir von der Kunst die Glaubhaftigkeit der Erscheinung, oder die Glaubhaftigkeit des Experiments? Das complex Natürliche, die kraftvolle Synthese, das Charakteristische mit seiner Verankerung im Unbegreiflichen — oder das angeordnet Willkürliche, die entkräftende Analyse, das Berechenbare und künstlich Vereinfachte? Die Kunst, die redend motivirt, ist von der zweiten Art, eine Schreibtisch- oder Laboratoriumskunst. Ein Versuch lässt sich einzeln durchsprechen; Reihenfolge und Zusammenwirken der Umstände lässt sich genau so vollständig darlegen, wie es der thatsächlichen Complication entspricht — warum? weil diese Complication einen sehr niederen Grad nicht überschreitet. Zweierlei, dreierlei ist vielleicht zu beachten, jedenfalls tritt nur eine discrete Anzahl von Gesetzlichkeiten ins Spiel. Die wirkliche Naturerscheinung, sei es Pflanze oder Thier, eine Wolkenbildung oder ein socialer Hergang, ist aber keine discrete Summe, sondern eine continuirliche, kein Zwei- oder Dreierlei, sondern ein Unendlich-Vielerlei, dessen einzelne Componenten nichtsdestoweniger vollkommen gesetzmässig verlaufen. Auf diese Gesetzmässigkeit im Einzelnen hat man auch in der Kunst den Ton gelegt und war dabei, wegen der Begrenztheit von Wort und Begriff, ans Discrete, Experimentelle gebunden, an das Bilden einer Summe aus sehr wenigen Factoren, an das Zerlegen eines Ganzen in sehr wenige, sehr grobe Theile. Die Menschen, denen man im Roman und auf der Bühne begegnet, sind Composita aus stark unterschiedlichen Antrieben und Charakteren, die einfach gruppenweise zusammengefasst und als nicht mehr auflösbar behandelt werden. Die Praxis der Chemie: irgend etwas spottet unserer Analyse, folglich ist es ein Element und bekommt einen Namen. So haben wir die bekannten Theatercharaktere und -temperamente, Herrschsucht, Güte, Düsterkeit, Leichtsinn, Geiz, Verschwendung — dass auch diese scheinbaren Elemente sehr complicirte Mischungen sind, dass die einfachste Handlung nur Abbreviatur eines unermesslich feinverzweigten Gegenspiels von Kräften und Widerständen ist, davon weiss der Theaterdichter nichts, und wenn er es weiss und eine höhere Complicationsstufe in Mensch und Ereigniss anstrebt, so bleibt er dennoch bei seiner Laboratoriumsmethode — er redet! Er fügt seiner Zwei und Drei die Vier hinzu, er will durchaus die Unbefangenheit und Natürlichkeit der Dinge, will sie aber auf constructivem Wege, er setzt Motiv auf Motiv, er spart nichts aus dem Notizbuch seiner piquanten Psychologismen — und doch steht am Ende seiner Bemühungen der errechnete Mensch, die combinirte Handlung, das unhistorische, untypische, uncharakteristische Geschehen. — Man ist unter Künstlern nicht dafür, mit der Kunst die Natur zu erklären; man will lieber Natur nachahmen, Natur sein. Macht man mit dieser Auffassung ernst, so verbanne man die scheinwissenschaftliche Kunst, die Paragraphen-Psychologie, die Kunst der redseligen Motivirung, der Schulcharakteristik, der breiten und simplificirenden Didaktik. Man verlange nicht, dass menschliche Bewegungen berechenbar seien wie Geschossbahnen, man bringe die Polyphonie des Geschehens nicht auf die primitive Formel einer Wirkung mit zwei oder drei Ursachen, oder eines Grundes mit zwei oder drei Folgen. Man trage kein Bedenken, die Natur als Natur zu geben, gewaltsam, sprunghaft, launisch, man wage die Kunstformen des Intermezzos, der Episode obenan zu stellen. Man suche nicht das Begreifliche, . das Einleuchtende, sondern das tyrannisch Thatsächliche, das Inconciliante, das nicht Deducible. Man lege den Schwerpunkt der Dinge dahin, wo ihn die Natur liegen hat: ins Unbewusste, ins Widersprüchliche, ins Spontane.
\par

